% LaTeX resume using res.cls
\documentclass[line,margin]{res}
\usepackage{hyperref}
\usepackage{enumitem}
\setlist[itemize]{topsep=-4pt, partopsep=0pt, itemsep=3pt, parsep=0pt}
\renewcommand{\familydefault}{\sfdefault}
\hypersetup{colorlinks=true, urlcolor=cyan}

\begin{document}
\name{Simon Crowe | Cloud-Native Software Engineer}
% \address can be used twice for additional contact info
\address{}
\address{\href{mailto://simon.r.crowe@protonmail.com}{simon.r.crowe@protonmail.com}}

\begin{resume}

\section{SUMMARY} Experienced in building and maintaining event-driven backends and platform tooling
                  on Kubernetes and AWS. While I'm comfortable writing product-facing APIs,
                  I work best at the intersection of application code and infrastructure:
                  designing processing pipelines, improving developer experience, and solving problems that
                  span service boundaries. I have a track record of raising engineering standards
                  through style guides, API conventions, and architectural patterns.
                  I feel at home with a DevOps culture and being on-call.

\section{EXPERIENCE} \textbf{Medtronic Digital Surgery} \\
                {\sl Senior Software Engineer} \hfill   Feb 2022 - present \\
                 \begin{itemize}
                   \item Designed and implemented a video processing orchestrator (Kubernetes, FastAPI, Kafka)
                         that replaced AWS MediaConvert in 97\% of cases. Saved ~\$280K/year, equivalent to 13\%
                         of the Surgical operating unit's FY26 cloud spend.
                   \item Instrumented the orchestrator with OpenTelemetry distributed tracing and structured logging
                         from the outset, enabling fault diagnosis without persistent job state.
                   \item Developed a fault-tolerant serverless solution for sending S3 objects to HTTP APIs as a
                         reusable Terraform module, adopted by multiple teams for event delivery to the data platform.
                   \item Re-architected a legacy Git/LFS assets server for multi-region cloud deployments.
                   \item Re-wrote legacy video processing integration tests in pytest, replacing an ad-hoc Python script.
                         Reduced worst-case failure feedback time from 90 to 30 minutes.
                   \item Authored a Python style guide (adopted as team standard)
                         and an engineering-wide document on HTTP API conventions.
                 \end{itemize}
                {\sl Platform Engineer} \hfill          Nov 2019 - Feb 2022 \\
                  \begin{itemize}
                  \item Led the implementation of MFA in line with OAuth 2.0 Best Current Practice.
                  \item Unified multiple video-sharing features, resolving data integrity issues
                        and simplifying the path for future sharing functionality.
                  \item Championed the service layer pattern to extract business logic from Django models,
                        improving testability and separation of concerns.
                  \end{itemize}

                \textbf{Verv} \\
                {\sl Python Developer} \hfill        Jun 2019 - Nov 2019 \\
                  \begin{itemize}
                  \item Python development: IoT energy monitoring device; Bluetooth GATT server PoC.
                  \end{itemize}

                \textbf{Ubiquity Press} \\
                {\sl Software Developer} \hfill        May 2018 - Jun 2019 \\
                  \begin{itemize}
                  \item Migrated legacy Django applications from manual SSH deployments
                        to continuous delivery pipelines using Docker, Kubernetes and Helm.
                  \item Built the foundations of an internal platform for deploying and managing
                        web apps.
                  \end{itemize}

	        \textbf{Titan Labs} \\
                {\sl Software Engineer} \hfill         Sep 2017 - Apr 2018 \\
                  \begin{itemize}
                  \item Independently implemented a PoC recommender system using LightFM.
                  \item Designed, built and documented REST APIs using Flask and OpenAPI 2.0.
                  \end{itemize}

\section{SKILLS} {\sl \textbf{Languages:}} Python, SQL, Bash\\
                 {\sl \textbf{Cloud \& Infrastructure:}} AWS, Kubernetes, Terraform, Docker, Helm\\
                 {\sl \textbf{Data \& Messaging:}} Kafka, PostgreSQL, Redis\\
                 {\sl \textbf{Frameworks \& Libraries:}} FastAPI, Django, pytest\\
                 {\sl \textbf{Observability:}} OpenTelemetry, structlog, OpenSearch, Cloudwatch, Honeycomb\\
                 {\sl \textbf{Some exposure to:}} Javascript, Go, Rust, GCP

\section{CERTIFICATIONS} CKAD: Certified Kubernetes Application Developer \hfill Feb 2025 \\
                         AWS Certified Solutions Architect -- Associate \hfill Mar 2024

\section{INTERESTS} Personal projects: \url{https://github.com/simoncrowe} \\
                    Active Blog: \url{https://simoncrowe.blog} \\
                    Cloud-focussed Engineering Blog: \url{https://simoncrowe.hashnode.dev} \\
                    Hobbies: running, strength training, reading, music, art

\section{EDUCATION} {\sl MA, Digital Culture} | Goldsmiths, University of London \hfill 2016 - 2017 \\
                    {\sl BA, Fine Art} | Loughborough University\hfill 2009 - 2013

\end{resume}
\end{document}
